% CUSTOM CODE FOR DRAWING NICE FIGURES
%
% (last updated: 5/1/2025)
%
% TRY THESE:
%
% ---- Meander picture for a set `S`.
%   \meanderS{4}{1/2, 3/2, 2/4}

%   \maxmeander{3/5}

%   \maxmeanderLower{3/5}


% ---- Layered graph for a set `S`
%   \gammaS{0/5/1, 2/2.5/2, 4/5/3, 2/0/4}{1/2, 3/2, 2/4}

%   \coloredgammaS{0/5/1/red, 2/2.5/2/red, 4/5/3/blue, 2/0/4/green}{1/2, 3/2, 2/4}


% ---- Seaweed shape, with test-functional highlighted
%   \Shape{1/2, 2/3}{1/1, 1/3, 2/1, 2/2, 3/3}{3/1, 3/2}


% ---- Seaweed shape, with eigenvalues highlighted
%   \eigenShape{1/2/\bgr1, 2/3/\bgr1, 1/1/\bl{0}, 1/3/2, 2/1/\rd{-1}, 2/2/\bl0, 3/3/\bl0}{3/1, 3/2}



%%%%%%%%%%%%%%%%%%%%%%%%%%%%%%%%%%%%%%%%%%
%%% some global definitions

% \usepackage{xstring}

\newcommand{\bgr}[1]{\textcolor{mygreen}{\ensuremath{\mathbf{#1}}}}
\newcommand{\gr}[1]{\textcolor{mygreen}{\ensuremath{#1}}}
\newcommand{\bl}[1]{\textcolor{myblue}{\ensuremath{#1}}}
\newcommand{\rd}[1]{\textcolor{red}{\ensuremath{#1}}}

\newcommand*{\getylength}[1]{
  \path let \p{y}=(0,1), \n{ylen}={veclen(\x{y},\y{y})}
    in \pgfextra{\xdef#1{\n{ylen}}};
}

%%%%%%%%%%%%%%%%%%%%%%%%%%%%%%%%%%%%%%%%%%
%%% Meander Macros
%
\newcommand{\meanderS}[3][.6]{%
\begin{tikzpicture}[scale=#1]
	\getylength{\mylen}
\tikzset{%
medge/.style = {-, line width=.5pt, black}, %mygreen!50!black,
vertex/.style = {line width=.5pt, font=\footnotesize, text=white, minimum size=0.2*\mylen, inner sep=0.05*\mylen}, %, mygreen!50!black
%shape=circle, draw, 
vertex text/.style = {font=\footnotesize, text=blue, minimum size=0.2*\mylen, inner sep=0.1*\mylen} %, mygreen!50!black
}
% vertices
\foreach \i in {1, ..., #2} {
	\node[vertex] (\i) at (\mylen*\i,0){8}; % invisible character of width "8"
	\node[vertex text] (n\i) at (\mylen*\i,0){$\mathsf\i$}; 
	}

%edges
\foreach \i/\j in {#3} {
    \pgfmathsetmacro{\x}{abs(\j-\i)}
    \draw[medge] (\i) to[bend left=40+5*\x] (\j); 
}
\end{tikzpicture}
}

\newcommand{\meanderD}[3][.6]{%
\begin{tikzpicture}[scale=#1]
	\getylength{\mylen}
\tikzset{%
medge/.style = {->, line width=.5pt, black}, %mygreen!50!black,
vertex/.style = {line width=.5pt, font=\footnotesize, text=white, minimum size=0.2*\mylen, inner sep=0.05*\mylen}, %, mygreen!50!black
%shape=circle, draw, 
vertex text/.style = {font=\footnotesize, text=blue, minimum size=0.2*\mylen, inner sep=0.1*\mylen} %, mygreen!50!black
}
% vertices
\foreach \i in {1, ..., #2} {
	\node[vertex] (\i) at (\mylen*\i,0){8}; % invisible character of width "8"
	\node[vertex text] (n\i) at (\mylen*\i,0){$\mathsf\i$}; 
	}

%edges
\foreach \i/\j in {#3} {
    \pgfmathsetmacro{\x}{abs(\j-\i)}
    \draw[medge] (\i) to[bend left=40+5*\x] (\j); 
}
\end{tikzpicture}
}

\newcommand{\meanderDUnderOverlay}[5][.6]{%
\begin{tikzpicture}[scale=#1]
	\getylength{\mylen}
\tikzset{%
medge/.style = {->, line width=.5pt, black}, %mygreen!50!black,
vertex/.style = {line width=.5pt, font=\footnotesize, text=white, minimum size=0.2*\mylen, inner sep=0.05*\mylen}, %, mygreen!50!black
%shape=circle, draw, 
vertex text/.style = {font=\footnotesize, text=blue, minimum size=0.2*\mylen, inner sep=0.1*\mylen} %, mygreen!50!black
}
{#4}
% vertices
\foreach \i in {1, ..., #2} {
	\node[vertex] (\i) at (\mylen*\i,0){8}; % invisible character of width "8"
	\node[vertex text] (n\i) at (\mylen*\i,0){$\mathsf\i$}; 
	}

%edges
\foreach \i/\j in {#3} {
    \pgfmathsetmacro{\x}{abs(\j-\i)}
    \draw[medge] (\i) to[bend left=40+5*\x] (\j); 
}
{#5}
\end{tikzpicture}
}


\newcommand{\maxmeander}[2][.6]{%
\begin{tikzpicture}[scale=#1]
	\getylength{\mylen}
\tikzset{%
medge/.style = {-, line width=.5pt, black}, %mygreen!50!black,
vertex/.style = {line width=.5pt, font=\footnotesize, text=white, minimum size=0.2*\mylen, inner sep=0.05*\mylen},
%, mygreen!50!black
%shape=circle, draw, 
vertex text/.style = {font=\footnotesize, text=blue, minimum size=0.2*\mylen, inner sep=0.1*\mylen} %, mygreen!50!black
}
\foreach \p/\q in {#2} {  % there should only be one!
    \pgfmathsetmacro{\n}{\p + \q}
    \pgfmathtruncatemacro{\nn}{floor(\n/2)}    
    \pgfmathtruncatemacro{\pp}{floor(\p/2)}    
    \pgfmathtruncatemacro{\qq}{floor(\q/2)}    
    % vertices
    \foreach \i in {1, ..., \n} {
       \node[vertex] (\i) at (\mylen*\i,0){8}; % invisible character of width "8"
	   \node[vertex text] (n\i) at (\mylen*\i,0){$\mathsf\i$}; 
	}

    %top edges
    \foreach \i [evaluate=\i as \j using int(\n+1-\i)] in {1, ..., \nn} {
        \draw[medge] (\i) to[bend left=75] (\j);
        % \draw[medge] (\i) to[bend left=50+3*(\j-\i)] (\j);
        %\draw[medge] (\i)  .. controls +(0.5, 0.5) and +(-0.5, 0.5) .. (\j);
    }

    %bottom edges, first block
    \foreach \i [evaluate=\i as \j using int(\p+1-\i)] in {1, ..., \pp} {
        \draw[medge] (\i) to[bend right=75] (\j); 
        % \draw[medge] (\i) to[bend right=50+3*(\j-\i)] (\j); 
    }

    %bottom edges, second block
    \foreach \i [evaluate=\i as \ii using int(\p+\i), 
                 evaluate=\i as \jj using int(\n+1-\i)] in {1, ..., \qq} {
        \draw[medge] (\ii) to[bend right=75] (\jj); 
        % \draw[medge] (\ii) to[bend right=50+3*(\jj-\ii)] (\jj); 
    }
    }
\end{tikzpicture}
}
\newcommand{\clipmaxmeander}[2][.6]{%
\begin{tikzpicture}[scale=#1]
	\getylength{\mylen}
\tikzset{%
medge/.style = {-, line width=.5pt, black}, %mygreen!50!black,
vertex/.style = {line width=.5pt, font=\footnotesize, text=white, minimum size=0.2*\mylen, inner sep=0.05*\mylen},
%, mygreen!50!black
%shape=circle, draw, 
vertex text/.style = {font=\footnotesize, text=blue, minimum size=0.2*\mylen, inner sep=0.1*\mylen} %, mygreen!50!black
}
\begin{scope}
    \clip(0,-5) rectangle (30,9);
\foreach \p/\q in {#2} {  % there should only be one!
    \pgfmathsetmacro{\n}{\p + \q}
    \pgfmathtruncatemacro{\nn}{floor(\n/2)}    
    \pgfmathtruncatemacro{\pp}{floor(\p/2)}    
    \pgfmathtruncatemacro{\qq}{floor(\q/2)}    
    % vertices
    \foreach \i in {1, ..., \n} {
       \node[vertex] (\i) at (\mylen*\i,0){8}; % invisible character of width "8"
	   \node[vertex text] (n\i) at (\mylen*\i,0){$\mathsf\i$}; 
	}

    %top edges
    \foreach \i [evaluate=\i as \j using int(\n+1-\i)] in {1, ..., \nn} {
        \draw[medge] (\i) to[bend left=70] (\j);
        % \draw[medge] (\i) to[bend left=50+3*(\j-\i)] (\j);
        %\draw[medge] (\i)  .. controls +(0.5, 0.5) and +(-0.5, 0.5) .. (\j);
    }

    %bottom edges, first block
    \foreach \i [evaluate=\i as \j using int(\p+1-\i)] in {1, ..., \pp} {
        \draw[medge] (\i) to[bend right=80] (\j); 
        % \draw[medge] (\i) to[bend right=50+3*(\j-\i)] (\j); 
    }

    %bottom edges, second block
    \foreach \i [evaluate=\i as \ii using int(\p+\i), 
                 evaluate=\i as \jj using int(\n+1-\i)] in {1, ..., \qq} {
        \draw[medge] (\ii) to[bend right=80] (\jj); 
        % \draw[medge] (\ii) to[bend right=50+3*(\jj-\ii)] (\jj); 
    }
    }
\end{scope}
\end{tikzpicture}
}


\newcommand{\maxmeanderD}[2][.6]{%
% `D` for directed edges
\begin{tikzpicture}[scale=#1]
	\getylength{\mylen}
\tikzset{%
medge/.style = {->, line width=.5pt, black}, %mygreen!50!black,
vertex/.style = {line width=.5pt, font=\footnotesize, text=white, minimum size=0.2*\mylen, inner sep=0.05*\mylen},
%, mygreen!50!black
%shape=circle, draw, 
vertex text/.style = {font=\footnotesize, text=blue, minimum size=0.2*\mylen, inner sep=0.1*\mylen} %, mygreen!50!black
}
\foreach \p/\q in {#2} {  % there should only be one!
    \pgfmathsetmacro{\n}{\p + \q}
    \pgfmathtruncatemacro{\nn}{floor(\n/2)}    
    \pgfmathtruncatemacro{\pp}{floor(\p/2)}    
    \pgfmathtruncatemacro{\qq}{floor(\q/2)}    
    % vertices
    \foreach \i in {1, ..., \n} {
       \node[vertex] (\i) at (\mylen*\i,0){8}; % invisible character of width "8"
	   \node[vertex text] (n\i) at (\mylen*\i,0){$\mathsf\i$}; 
	}

    %top edges
    \foreach \i [evaluate=\i as \j using int(\n+1-\i)] in {1, ..., \nn} {
        \draw[medge] (\i) to[bend left=75] (\j);
        % \draw[medge] (\i) to[bend left=50+3*(\j-\i)] (\j);
        %\draw[medge] (\i)  .. controls +(0.5, 0.5) and +(-0.5, 0.5) .. (\j);
    }

    %bottom edges, first block
    \foreach \i [evaluate=\i as \j using int(\p+1-\i)] in {1, ..., \pp} {
        \draw[medge] (\j) to[bend left=75] (\i); 
        % \draw[medge] (\j) to[bend left=50+3*(\j-\i)] (\i); 
    }

    %bottom edges, second block
    \foreach \i [evaluate=\i as \ii using int(\p+\i), 
                 evaluate=\i as \jj using int(\n+1-\i)] in {1, ..., \qq} {
        \draw[medge] (\jj) to[bend left=75] (\ii); 
        % \draw[medge] (\jj) to[bend left=50+3*(\jj-\ii)] (\ii); 
    }
    }
    \end{tikzpicture}
}



\newcommand{\maxmeanderLower}[2][.6]{%
\begin{tikzpicture}[scale=#1]
	\getylength{\mylen}
\tikzset{%
medge/.style = {-, line width=.5pt, black}, %mygreen!50!black,
vertex/.style = {line width=.5pt, font=\footnotesize, text=white, minimum size=0.2*\mylen, inner sep=0.05*\mylen},
%, mygreen!50!black
%shape=circle, draw, 
vertex text/.style = {font=\footnotesize, text=blue, minimum size=0.2*\mylen, inner sep=0.1*\mylen} %, mygreen!50!black
}
\foreach \p/\q in {#2} {  % there should only be one!
    \pgfmathsetmacro{\n}{\p + \q}
    \pgfmathtruncatemacro{\nn}{floor(\n/2)}    
    \pgfmathtruncatemacro{\pp}{floor(\p/2)}    
    \pgfmathtruncatemacro{\qq}{floor(\q/2)}    
    % vertices
    \foreach \i in {1, ..., \n} {
       \node[vertex] (\i) at (\mylen*\i,0){8}; % invisible character of width "8"
	   \node[vertex text] (n\i) at (\mylen*\i,0){$\mathsf\i$}; 
	}

    %top edges
    \foreach \i [evaluate=\i as \j using int(\n+1-\i)] in {1, ..., \nn} {
        \draw[medge] (\i) to[bend right=75] (\j);
        % \draw[medge] (\i) to[bend right=50+3*(\j-\i)] (\j);
        %\draw[medge] (\i)  .. controls +(0.5, 0.5) and +(-0.5, 0.5) .. (\j);
    }

    %bottom edges, first block
    \foreach \i [evaluate=\i as \j using int(\p+1-\i)] in {1, ..., \pp} {
        \draw[medge] (\i) to[bend left=75] (\j); 
        % \draw[medge] (\i) to[bend left=50+3*(\j-\i)] (\j); 
    }

    %bottom edges, second block
    \foreach \i [evaluate=\i as \ii using int(\p+\i), 
                 evaluate=\i as \jj using int(\n+1-\i)] in {1, ..., \qq} {
        \draw[medge] (\ii) to[bend left=75] (\jj); 
        % \draw[medge] (\ii) to[bend left=50+3*(\jj-\ii)] (\jj); 
    }
    }
    \end{tikzpicture}
}


\newcommand{\MaxmeanderLower}[3][.6]{%
\begin{tikzpicture}[scale=#1]
	\getylength{\mylen}
\tikzset{%
medge/.style = {-, line width=.5pt, black}, %mygreen!50!black,
vertex/.style = {line width=.5pt, font=\footnotesize, text=white, minimum size=0.1*\mylen, inner sep=0.02*\mylen},
%, mygreen!50!black
%shape=circle, draw, 
vertex text/.style = {font=\footnotesize, text=blue, minimum size=0.2*\mylen, inner sep=0.1*\mylen} %, mygreen!50!black
}
\foreach \p/\q in {#2} {  % there should only be one!
    \pgfmathtruncatemacro{\n}{\p + \q}
    \pgfmathtruncatemacro{\px}{\p + 1}    
    \pgfmathtruncatemacro{\qx}{\q + 1}    
    \pgfmathtruncatemacro{\nn}{floor(\n/2)}    
    \pgfmathtruncatemacro{\pp}{floor(\p/2)}    
    \pgfmathtruncatemacro{\qq}{floor(\q/2)}    
    % vertices
    \ifthenelse{#3=0}{%
    % bullet the first entries
        \foreach \i in {1, ..., \p} {
           \node[vertex] (\i) at (\mylen*\i,0){o}; % invisible character of width "o"
    	   \node[vertex text] (n\i) at (\mylen*\i,0){\!${}_{}\bs\cdot$}; 
    	}
        \foreach \i in {1, ..., \q} {
           \pgfmathtruncatemacro\ix{\p + \i}
           \node[vertex] (\ix) at (\mylen*\ix,0){o}; % invisible character of width "o"
    	   \node[vertex text] (n\ix) at (\mylen*\ix,0){\!$\mathsf\i$}; 
    	}
    }{%
    % bullet the last entries
        \foreach \i in {1, ..., \p} {
           \node[vertex] (\i) at (\mylen*\i,0){o}; % invisible character of width "o"
    	   \node[vertex text] (n\i) at (\mylen*\i,0){$\mathsf\i$\!}; 
    	}
        \foreach \i in {\px, ..., \n} {
           \node[vertex] (\i) at (\mylen*\i,0){o}; % invisible character of width "o"
    	   \node[vertex text] (n\i) at (\mylen*\i,0){$\bs\cdot{}_{}$\!}; 
    	}
    }
    %top edges
    \foreach \i [evaluate=\i as \j using int(\n+1-\i)] in {1, ..., \nn} {
        \draw[medge] (\i) to[bend right=75] (\j);
        % \draw[medge] (\i) to[bend right=65+3*(\j-\i)] (\j);
        %\draw[medge] (\i)  .. controls +(0.5, 0.5) and +(-0.5, 0.5) .. (\j);
    }

    %bottom edges, first block
    \foreach \i [evaluate=\i as \j using int(\p+1-\i)] in {1, ..., \pp} {
        % \draw[medge] (\i) to[in=90,out=90] (\j); 
        \draw[medge] (\i) to[bend left=75] (\j); 
        % \draw[medge] (\i) to[bend left=65+3*(\j-\i)] (\j); 
    }

    %bottom edges, second block
    \foreach \i [evaluate=\i as \ii using int(\p+\i), 
                 evaluate=\i as \jj using int(\n+1-\i)] in {1, ..., \qq} {
        % \draw[medge] (\ii) to[in=90,out=90] (\jj); 
        \draw[medge] (\ii) to[bend left=75] (\jj); 
        % \draw[medge] (\ii) to[bend left=65+3*(\jj-\ii)] (\jj); 
    }
    }
    \end{tikzpicture}
}


\newcommand{\maxmeanderRetracting}[5][.6]{%
\begin{tikzpicture}[scale=#1]
  \getylength{\mylen}
  \tikzset{%
    medge/.style  = {-, line width=.5pt, black},
    vertex/.style = {line width=.5pt, font=#5, text=white,
                     minimum size=0.2*\mylen, inner sep=0.05*\mylen},
    vertex text/.style = {font=#5, text=blue,
                     minimum size=0.1*\mylen, inner sep=0.1*\mylen}
  }
  \foreach \p/\q in {#2} {
    \pgfmathsetmacro{\n}{\p + \q}
    \pgfmathtruncatemacro{\nn}{floor(\n/2)}
    \pgfmathtruncatemacro{\pp}{floor(\p/2)}
    \pgfmathtruncatemacro{\qq}{floor(\q/2)}

    % Strip units from \mylen for coordinate arithmetic
    \pgfmathsetmacro{\mylennum}{\mylen/1pt}

    % ---- Baseline nodes 1 .. p ----
    \foreach \i in {1, ..., \p} {
      \node[vertex]      (\i)  at (\mylennum*\i pt, 0) {8};
      \node[vertex text] (n\i) at (\mylennum*\i pt, 0) {$\mathsf{\i}$};
    }
    
    % ---- former baseline nodes p+1 .. p+q ----
    \foreach \j in {1, ..., \q} {
      \pgfmathtruncatemacro{\pj}{\p + \j}
      \node[vertex text] (f\j) at (\mylennum*\pj pt, 0) {\color{blue!40}{$\mathsf{\pj}$}};
    }

    % ---- Tongue nodes p+1 .. n, on 15-deg ray from node nn+1 ----
    \pgfmathtruncatemacro{\apex}{ceil(\n/2)}
    \foreach \k in {1, ..., \q} {
      \pgfmathtruncatemacro{\idx}{\p + \k}
      \pgfmathsetmacro{\dist}{\idx - \apex}
      \pgfmathsetmacro{\tx}{\mylennum*\apex + (\dist-#4)*\mylennum*cos(#3)}
      \pgfmathsetmacro{\ty}{(\dist-#4)*\mylennum*sin(#3)}
      \node[vertex]      (\idx) at (\tx pt, \ty pt) {8};
      \node[vertex text] (n\idx) at (\tx pt, \ty pt) {\large$\boldsymbol{\cdot}$};
    }

    % ---- Top arcs: pair i with n+1-i for i=1..\nn ----
    \foreach \i [evaluate=\i as \j using int(\n+1-\i)] in {1, ..., \nn} {
      \draw[medge] (\i) to[bend left=40+2*(\j-\i)] (\j);
    }

    % ---- Bottom arcs, left block: pair i with p+1-i ----
    \foreach \i [evaluate=\i as \j using int(\p+1-\i)] in {1, ..., \pp} {
      \draw[medge] (\i) to[bend right=50+3*(\j-\i)] (\j);
    }

    % ---- Bottom arcs, tongue block: pair p+k with p+q+1-k ----
    \foreach \k in {1, ..., \qq} {
      \pgfmathtruncatemacro{\ia}{\p + \k}
      \pgfmathtruncatemacro{\ib}{\p + \q + 1 - \k}
      \pgfmathtruncatemacro{\bendamt}{50+3*(\ib-\ia)}
      \draw[medge] (\ia) to[bend right=\bendamt] (\ib);
    }
  }
\end{tikzpicture}
}
% try: \maxmeanderRetracting[.35]{7/4}{35}{.4}{\tiny}

%%%%%%%%%%%%%%%%%%%%%%%%%%%%%%%%%%%%%%%%%%
%%% LayeredGraph Macros
%
\newcommand{\gammaS}[3][.6]{%
\raisebox{-.5\height}{%
\begin{tikzpicture}[scale=#1]
	\getylength{\mylen}
\tikzset{%
medge/.style = {->, >=latex', line width=.5pt, black}, %mygreen!50!black,
vertex/.style = {shape=circle, draw, line width=.5pt, font=\footnotesize, text=white, minimum size=0.2*\mylen, inner sep=0.1*\mylen}, %, mygreen!50!black
vertex text/.style = {font=\footnotesize, text=black, minimum size=0.2*\mylen, inner sep=0.1*\mylen} %, mygreen!50!black
}
% vertices
\foreach \i/\j/\x in {#2} {
	\node[vertex] (\x) at (\i,\j){8}; % invisible character of width "8"
	\node[vertex text] (n\x) at (\i,\j){$\mathsf\x$}; % use two nodes to better control vertex size
	}
%NOTE: this would yield a hidden node (since no "draw")
% \node [shape=circle,minimum size=1.5em] (h) at (1,1.5) {}; 

%edges
\foreach \x/\y in {#3} {
	\draw[medge] (\x) to (\y); % or you could bend it with `to[bend left]`
	}
\end{tikzpicture}
}}
% try: \gammaS{0/5/1, 2/2.5/2, 4/5/3, 2/0/4}{1/2, 3/2, 2/4}


\newcommand{\gammaSUnderOverlay}[5][.6]{%
\raisebox{-.5\height}{%
\begin{tikzpicture}[scale=#1]
	\getylength{\mylen}
\tikzset{%
medge/.style = {->, > = latex', line width=.5pt, black}, %mygreen!50!black,
vertex/.style = {shape=circle, draw, fill=white,line width=.5pt, font=\footnotesize, text=white, minimum size=0.2*\mylen, inner sep=0.1*\mylen}, %, mygreen!50!black
vertex text/.style = {font=\footnotesize, text=black, minimum size=0.2*\mylen, inner sep=0.1*\mylen} %, mygreen!50!black
}
% underlay
#4
% vertices
\foreach \i/\j/\x in {#2} {
	\node[vertex] (\x) at (\i,\j){8}; % invisible character of width "8"
	\node[vertex text] (n\x) at (\i,\j){$\mathsf\x$}; % use two nodes to better control vertex size
	}
%NOTE: this would yield a hidden node (since no "draw")
% \node [shape=circle,minimum size=1.5em] (h) at (1,1.5) {}; 

%edges
\foreach \x/\y in {#3} {
	\draw[medge] (\x) to (\y); % or you could bend it with `to[bend left]`
	}
% overlay
#5
\end{tikzpicture}
}}
% try: \gammaS{0/5/1, 2/2.5/2, 4/5/3, 2/0/4}{1/2, 3/2, 2/4}

\newcommand{\coloredgammaS}[3][.6]{%
\raisebox{-.5\height}{%
\begin{tikzpicture}[scale=#1]
	\getylength{\mylen}
\tikzset{%
medge/.style = {->, > = latex', line width=.5pt, black}, %mygreen!50!black,
vertex/.style = {shape=circle, draw, line width=.5pt, font=\footnotesize, minimum size=0.2*\mylen, inner sep=0.1*\mylen}, %, mygreen!50!black
vertex text/.style = {font=\footnotesize, text=black, minimum size=0.2*\mylen, inner sep=0.1*\mylen} %, mygreen!50!black
}
% vertices
\foreach \i/\j/\x\c in {#2} {
	\node[vertex,fill=\c, text=\c] (\x) at (\i,\j){8}; % invisible character of width "8"
	\node[vertex text] (n\x) at (\i,\j){$\mathsf\x$}; % use two nodes to better control vertex size
	}
%NOTE: this would yield a hidden node (since no "draw")
% \node [shape=circle,minimum size=1.5em] (h) at (1,1.5) {}; 

%edges
\foreach \x/\y in {#3} {
	\draw[medge] (\x) to (\y); % or you could bend it with `to[bend left]`
	}
\end{tikzpicture}
}}
% try: \coloredgammaS{0/5/1/red, 2/2.5/2/red, 4/5/3/blue, 2/0/4/green}{1/2, 3/2, 2/4}

\newcommand{\coloredgammaSUnderlay}[4][.6]{%
\raisebox{-.5\height}{%
\begin{tikzpicture}[scale=#1]
	\getylength{\mylen}
\tikzset{%
medge/.style = {->, > = latex', line width=.5pt, black}, %mygreen!50!black,
vertex/.style = {shape=circle, draw, line width=.5pt, font=\footnotesize, minimum size=0.2*\mylen, inner sep=0.1*\mylen}, %, mygreen!50!black
vertex text/.style = {font=\footnotesize, text=black, minimum size=0.2*\mylen, inner sep=0.1*\mylen} %, mygreen!50!black
}
% vertices
\foreach \i/\j/\x\c in {#2} {
	\node[vertex,fill=\c, text=\c] (\x) at (\i,\j){8}; % invisible character of width "8"
	\node[vertex text] (n\x) at (\i,\j){$\mathsf\x$}; % use two nodes to better control vertex size
	}
%NOTE: this would yield a hidden node (since no "draw")
% \node [shape=circle,minimum size=1.5em] (h) at (1,1.5) {}; 

{#4}
%edges
\foreach \x/\y in {#3} {
	\draw[medge] (\x) to (\y); % or you could bend it with `to[bend left]`
	}
\end{tikzpicture}
}}
% try: \coloredgammaSUnderlay{0/5/1/red, 2/2.5/2/red, 4/5/3/blue, 2/0/4/green}{1/2, 3/2, 2/4}{\draw (2) -- (3);}


\newcommand{\orderideals}[4][.6]{%
\raisebox{-.5\height}{%
\begin{tikzpicture}[scale=#1]
\tikzset{%
hedge/.style = {-, line width=.5pt, black}, %mygreen!50!black,
vertex/.style = {black, font=\small, fill=white, inner sep=2pt}, %draw=red, line width=.1pt, 
}
% vertices
\foreach \i/\j/\x/\l in {#3} {
    \StrSubstitute{\l}{,}{{,}}[\ltight]%
    \if#20
        \node[vertex] (\x) at (\i,\j){$\{\mathsf{\l}\}$};
    \else
        \node[vertex] (\x) at (\i,\j){$\mathsf{\ltight}$};
    \fi
    }
%edges
\foreach \x/\y in {#4} {
	\draw[hedge] (\x) to (\y); 
	}
\end{tikzpicture}
}}
% try changing the second arugment (0 or 1):
% =0: 
	% \orderideals[.45]{0}{-2/4/12/{1,2}, 0/6/123/{1,2,3}, 0/2/2/{2}, 2/4/23/{2,3}, 0/0/e/{\emptyset}}{2/e, 12/2, 23/2, 123/12, 123/23}
% =1: 	
	% \orderideals[.45]{1}{-2/4/12/{1,2}, 0/6/123/{1,2,3}, 0/2/2/{2}, 2/4/23/{2,3}, 0/0/e/{\emptyset}}{2/e, 12/2, 23/2, 123/12, 123/23}


\newcommand{\coloredorderideals}[4][.6]{%
\raisebox{-.5\height}{%
\begin{tikzpicture}[scale=#1]
\tikzset{%
hedge/.style = {-, line width=.5pt, black}, %mygreen!50!black,
vertex-bg/.style = {white, shape=circle, font=\small, fill=white, inner sep=1pt}, vertex/.style = {black, font=\small, inner sep=1pt}, %draw=red, line width=.1pt, 
}
% vertices
\foreach \i/\j/\x/\l/\c in {#3} {
   \StrSubstitute{\l}{,}{{,}}[\ltight]%
    \if#20
        \node[vertex-bg, shape=circle,fill=\c] (bg\x) at (\i,\j){\phantom{8}};
        \node[vertex] (\x) at (\i,\j){$\{\mathsf{\l}\}$};
    \else
        \node[vertex-bg, shape=circle,fill=\c] (bg\x) at (\i,\j){\phantom{8}};
        \node[vertex] (\x) at (\i,\j){$\mathsf{\ltight}$};
    \fi
    }
%edges
\foreach \x/\y in {#4} {
	\draw[hedge] (\x) to (\y); 
	}
\end{tikzpicture}
}}
% try: \coloredorderideals[.45]{1}{-2/4/12/{1,2}/white, 0/6/123/{1,2,3}/white, 0/2/2/{2}/blue!20, 2/4/23/{2,3}/white, 0/0/e/{\emptyset}/blue!20}{2/e, 12/2, 23/2, 123/12, 123/23}



%%%%%%%%%%%%%%%%%%%%%%%%%%%%%%%%%%%%%%%%%%
%%% Shape Macros
%
\newcommand{\eigenShape}[3][.7]{%
\raisebox{-.5\height}{%
\begin{tikzpicture}[scale=#1]
	\getylength{\mylen}
\tikzset{%
blank/.style = {shape=rectangle, font=\tiny, text=gray!80, minimum size=0.5*\mylen, inner sep=0.1*\mylen},
box/.style = {shape=rectangle, line width=.5pt, white, fill=gray!20, draw},
box text/.style = {font=\footnotesize, text=black, minimum size=0.2*\mylen, inner sep=0*\mylen}
%box text/.style = {font=\normalsize, text=black, minimum size=0.2*\mylen, inner sep=0*\mylen, scale=#1}
}
\foreach \i/\j/\x in {#2} {
	\draw[box] (\j-.5,-\i-.5) rectangle ++(1,1);
	\node[box text] (\i\j) at (\j,-\i){$\x$};
	}
\foreach \i/\j in {#3}
	\node[blank] (\i\j) at (\j,-\i){$\cdot$};
\end{tikzpicture}
}}
% try: \eigenShape{1/2/\bgr1, 2/3/\bgr1, 1/1/\bl{0}, 1/3/2, 2/1/\rd{-1}, 2/2/\bl0, 3/3/\bl0}{3/1, 3/2}

\newcommand{\eigenShapeOverlay}[4][.7]{%
\raisebox{-.5\height}{%
\begin{tikzpicture}[scale=#1]
	\getylength{\mylen}
\tikzset{%
blank/.style = {shape=rectangle, font=\tiny, text=gray!80, minimum size=0.5*\mylen, inner sep=0.1*\mylen},
box/.style = {shape=rectangle, line width=.5pt, white, fill=gray!20, draw},
box text/.style = {font=\footnotesize, text=black, minimum size=0.2*\mylen, inner sep=0*\mylen}
%box text/.style = {font=\normalsize, text=black, minimum size=0.2*\mylen, inner sep=0*\mylen, scale=#1}
}
\foreach \i/\j/\x in {#2} {
	\draw[box] (\j-.5,-\i-.5) rectangle ++(1,1);
	\node[box text] (\i\j) at (\j,-\i){$\x$};
	}
\foreach \i/\j in {#3}
	\node[blank] (\i\j) at (\j,-\i){$\cdot$};
% overlay
#4
\end{tikzpicture}
}}
% try: \eigenShapeOverlay{1/2/\bgr1, 2/3/\bgr1, 1/1/\bl{0}, 1/3/2, 2/1/\rd{-1}, 2/2/\bl0, 3/3/\bl0}{3/1, 3/2}{\fill[magenta!90, nearly transparent] (7,-6) circle (.55);}

\newcommand{\Shape}[4][.7]{%
\begin{tikzpicture}[scale=#1]
	\getylength{\mylen}
\tikzset{%
blank/.style = {shape=rectangle, font=\tiny, text=gray!80, minimum size=0.5*\mylen, inner sep=0.1*\mylen},
box/.style = {shape=rectangle, line width=.5pt, white, fill=gray!20, draw},
box text/.style = {font=\footnotesize, text=black, minimum size=0.2*\mylen, inner sep=0*\mylen}
%box text/.style = {font=\normalsize, text=black, minimum size=0.2*\mylen, inner sep=0*\mylen, scale=#1}
}
\foreach \i/\j\x in {#2} {
	\draw[box] (\j-.5,-\i-.5) rectangle ++(1,1);
	\node[box text] (\i\j) at (\j,-\i){\gr{\i{\scriptscriptstyle\,}\j}};
	}
\foreach \i/\j in {#3} {
	\draw[box] (\j-.5,-\i-.5) rectangle ++(1,1);
	\ifthenelse{\i=\j}{\node[box text] (\i\j) at (\j,-\i){\bl{\times}};}{\node[box text] (\i\j) at (\j,-\i){$\times$};}
	}
\foreach \i/\j in {#4}
	\node[blank] (\i\j) at (\j,-\i){$\cdot$};
\end{tikzpicture}
}
% try: \Shape{1/2, 2/3}{1/1, 1/3, 2/1, 2/2, 3/3}{3/1, 3/2}

\newcommand{\ShapeOverlay}[5][.7]{%
\begin{tikzpicture}[scale=#1]
	\getylength{\mylen}
\tikzset{%
blank/.style = {shape=rectangle, font=\tiny, text=gray!80, minimum size=0.5*\mylen, inner sep=0.1*\mylen},
box/.style = {shape=rectangle, line width=.5pt, white, fill=gray!20, draw},
box text/.style = {font=\footnotesize, text=black, minimum size=0.2*\mylen, inner sep=0*\mylen}
%box text/.style = {font=\normalsize, text=black, minimum size=0.2*\mylen, inner sep=0*\mylen, scale=#1}
}
\foreach \i/\j\x in {#2} {
	\draw[box] (\j-.5,-\i-.5) rectangle ++(1,1);
	\node[box text] (\i\j) at (\j,-\i){\gr{\i{\scriptscriptstyle\,}\j}};
	}
\foreach \i/\j in {#3} {
	\draw[box] (\j-.5,-\i-.5) rectangle ++(1,1);
	\ifthenelse{\i=\j}{\node[box text] (\i\j) at (\j,-\i){\bl{\times}};}{\node[box text] (\i\j) at (\j,-\i){$\times$};}
	}
\foreach \i/\j in {#4}
	\node[blank] (\i\j) at (\j,-\i){$\cdot$};
%% overlay stuff
#5
\end{tikzpicture}
}
%%%%%%%%%%%%%%%%%%%%%%%%%%%%%%%%%%%%%%%%%%



%%%%%%%%%%%%%%%%%%%%%%%%%%%%%%%%%%%%%%%%%%
%%% Fence Poset Macros
%
\newcommand{\fenceposet}[2][.6]{%
\raisebox{.5\height}{%
\begin{tikzpicture}[scale=#1]
\tikzset{%
fedge/.style = {-, line width=.5pt, black}, %mygreen!50!black,
vertex/.style = {blue, font=\small, fill=white, inner sep=1pt}, %draw=red, line width=.1pt, 
}
    \draw[fedge] (0,0)
    \foreach \I in {#2} {
      \ifdim\I pt > 0pt -- ++(1,1) \else -- ++(1,-1) \fi
    };
    \node[vertex] at (0,0) {$\mathsf1$};
    \coordinate (current) at (0,0);
    \foreach[count=\t from 2] \I in {#2} {
      \ifdim\I pt > 0pt
        \coordinate (current) at ($(current) + (1,1)$);
      \else
        \coordinate (current) at ($(current) + (1,-1)$);
      \fi
      \node[vertex] at (current) {$\mathsf\t$};
    }
  \end{tikzpicture}%
}}
% try: \fenceposet{1,-1,-1,1,1,1,-1}


\newcommand{\fenceposetExtended}[3][.6]{%
\raisebox{.5\height}{%
\begin{tikzpicture}[scale=#1]
\tikzset{%
  fedge/.style   = {-, line width=.5pt, black},
  vertex/.style  = {blue, font=\tiny, fill=white, inner sep=1pt},
  fedgeX/.style  = {-, line width=.6pt, magenta!90!black},
  vertexX/.style = {magenta!90!black, font=\tiny, fill=white, inner sep=.5pt},
}
  % First pass: build coordinates for main path (#2)
  \coordinate (F1) at (0,0);
  \foreach \I [count=\k from 2] in {#2} {
    \pgfmathtruncatemacro{\kprev}{\k-1}
    \ifdim\I pt > 0pt
      \coordinate (F\k) at ($(F\kprev)+(1,1)$);
    \else
      \coordinate (F\k) at ($(F\kprev)+(1,-1)$);
    \fi
    \global\edef\maincount{\k}%
  }
  % One extra step for the extended edge (#3)
  \pgfmathtruncatemacro{\extk}{\maincount+1}
  \ifdim#3pt > 0pt
    \coordinate (F\extk) at ($(F\maincount)+(1,1)$);
  \else
    \coordinate (F\extk) at ($(F\maincount)+(1,-1)$);
  \fi
  % Draw main edges
  \draw[fedge] (F1)
    \foreach \k in {2,...,\maincount} { -- (F\k) };
  % Draw extended edge
  \draw[fedgeX] (F\maincount) -- (F\extk);
  % Draw main nodes
  \foreach \k in {1,...,\maincount} {
    \node[vertex] at (F\k) {$\mathsf{\k}$};
  }
  % Draw extended node
  % \node[vertexX,shape=circle] at (F\extk) {$\boldsymbol{\times}$};
  \node[vertexX,shape=circle] at (F\extk) {$\mathsf{\extk}$};
\end{tikzpicture}%
}}% try: \fenceposetExtended{1,-1,-1,1,1,1}{-1}



%%% code for pyramids and such
%\usepackage{xparse}

\def\skylineboxsize{0.475}   % cm ? side length of each unit box
\def\skylinegap{0.125}       % cm ? gap between adjacent boxes

\newcommand{\skyline}[1]{%
  \begin{tikzpicture}[
      x=1cm, y=1cm,
      box/.style={
        draw=black, line width=0.6pt,
        fill=white,
        minimum size=\skylineboxsize cm,
        inner sep=0pt,
      },
      baseline=(current bounding box.south),
  ]
    \foreach \Y [count=\xi from 0] in {#1} {%
      \foreach \yi in {1,...,\Y} {%
        \pgfmathsetmacro{\bx}{\xi*(\skylineboxsize+\skylinegap)}%
        \pgfmathsetmacro{\by}{(\yi-1)*(\skylineboxsize+\skylinegap)}%
        \node[box] at (\bx,\by) {};%
      }%
    }%
  \end{tikzpicture}%
}

\newcommand{\fskyline}[1]{%
  \begin{tikzpicture}[
      x=1cm, y=1cm,
      box/.style={
        draw=black, line width=0.6pt,
        fill=white,
        minimum size=\skylineboxsize cm,
        inner sep=0pt,
      },
      baseline=(current bounding box.south),
  ]
    \foreach \Labels [count=\xi from 0] in {#1} {%
      \foreach \lbl [count=\yi from 1] in \Labels {%
        \pgfmathsetmacro{\bx}{\xi*(\skylineboxsize+\skylinegap)}%
        \pgfmathsetmacro{\by}{(\yi-1)*(\skylineboxsize+\skylinegap)}%
        \node[box] (\lbl) at (\bx,\by) {\small\lbl};%
      }%
    }%
  \end{tikzpicture}%
}

\newcommand{\fcskyline}[1]{%
  \begin{tikzpicture}[
      x=1cm, y=1cm,
      box/.style={draw=black, line width=0.6pt,
                  minimum size=\skylineboxsize cm, inner sep=0pt},
      baseline=(current bounding box.south),
  ]
    \foreach \Pairs [count=\xi from 0] in {#1} {%
      \foreach \lbl/\col [count=\yi from 1] in \Pairs {%
        \pgfmathsetmacro{\bx}{\xi*(\skylineboxsize+\skylinegap)}%
        \pgfmathsetmacro{\by}{(\yi-1)*(\skylineboxsize+\skylinegap)}%
        \node[box, fill=\col!30, name=n-\xi-\yi] at (\bx,\by) {\small\lbl};
        }%
    }%
  \end{tikzpicture}%
}

\newcommand{\fcskylineOverlay}[2]{%
  \begin{tikzpicture}[
      x=1cm, y=1cm,
      box/.style={draw=black, line width=0.6pt,
                  minimum size=\skylineboxsize cm, inner sep=0pt},
      baseline=(current bounding box.south),
  ]
    \foreach \Pairs [count=\xi from 0] in {#1} {%
      \foreach \lbl/\col [count=\yi from 1] in \Pairs {%
        \pgfmathsetmacro{\bx}{\xi*(\skylineboxsize+\skylinegap)}%
        \pgfmathsetmacro{\by}{(\yi-1)*(\skylineboxsize+\skylinegap)}%
        \node[box, fill=\col!30, name=n-\xi-\yi] at (\bx,\by) {\small\lbl};
      }%
    }%
    #2
  \end{tikzpicture}%
}




%%%% HOW TO DRAW A TREE
%
%\begin{tikzpicture}
%\node {root}
%child {
%node {left}
%edge from parent
%node[left] {a}
%node[right] {b}
%}
%child {
%node {right}
%child {
%node {child}
%edge from parent
%node[left] {c}
%}
%child {node {child}}
%edge from parent
%node[near end] {x}
%};
%\end{tikzpicture}

